\chapterNote{Based on a talk given on March 28, 2019}

\chapter{Know, Watch, Let Go}
\tocChapterNote{Based on a talk given on March 28, 2019}
%\tocChapterNote{March 28, 2019}


This is the last Lunar Observance Day of our Winter Retreat, and we're about to transition to moving back into more engagement---guests, visitors, work routines, the ``normal'' schedule of the monastery. We can watch the energy as it starts to build momentum towards that transition. Depending on the individual, there may have been a fair amount of settling, quietude, and gathering of energy during retreat, and once the restraints of noble silence and a reduced schedule are relaxed, the energy can start to bubble over.

I think we've all experienced that on retreats. When the noble silence is lifted and we're resuming normal activities, we find ourselves talking loud and fast, blurting out lots of runaway-train ideas. Watch this transition; see how it feels and be mindful.

There is a famous line from Ajahn Sucitto, just before a Winter Retreat started at Chithurst Monastery. Somebody asked him, probably a week or two before the retreat was supposed to ensue: ``Ajahn, when is the retreat going to start? What day is the retreat actually going to begin?'' His reply was: ``For some people, the retreat has already started. For some people, it will never begin.''

We could probably say the same thing for the ending of the retreat. For some people, the retreat has already ended, or maybe never started. For others, it will continue on. That, in my mind, is the goal: to keep some of the retreat feeling. We can protect that feeling without clinging to it and without lamenting any change that does occur.

I think most people here had a fruitful retreat, if not always easy. The community as a whole was committed to the practice and was definitely on retreat. We can work to maintain any skillful qualities that we've learned and to create habits to nurture and protect them. One way to do that is to notice the difference between outward-going momentum and the stillness that has been nurtured through various ways in the past three months.

During one of the readings a few days ago, there was a question about what the term ``\emph{saṅkhāra}'' refers to. Also, the word ``\emph{saṅkhāra}'' appears in the funeral chanting in the steps of dependent origination. It's the second factor: based on ignorance, \emph{saṅkhāras} arise. It also appears in the last reflection on impermanence that we do as part of the funeral chanting: \emph{aniccā vata saṅkhārā}. It's a word that is used broadly, variously defined as conditions, fabrications, mind conditioners, volitional formations, and mental formations. It's used differently in different circumstances, but it has the quality of flow to it.

Whether it's used as a general reference to all conditioned phenomena, or to more specific phenomena, it has the quality of flow, of movement, rolling on. It's very useful to view it as a fluid experience rather than trying to pin it down as a static object. That can be difficult to do. We often speak in terms of nouns: \emph{Aniccā vata saṅkhārā}, ``all compounded things are impermanent.'' \emph{Saṅkhāras}---compounded things, creations, mind states---those nouns tend to solidify our perception of saṅkhāras.

All of these mental phenomena have the quality of movement and flow. Speaking from the perspective of experience, there is nothing that's solid. It's all an experience of movement and flow. The quality of \emph{cetanā}---intention, volition---is one of saṅkhāra's primary characteristics. Intention, the mental impulse or driving force behind all action, feeds the saṅkhāra process.

Saṅkhāra, in its active phase, produces mental qualities and activity. It serves as the basis for action that manifests through body, speech, or mind.

Sometimes saṅkhāras are experienced as \emph{vipāka}, resultant conditions. Perceptions and feelings are the end result of the intentional process. They are resultant conditions, but are still part of the saṅkhāra flow. This rolling on of productive and resultant experience conditions even further intentional action.

If the process is driven by greed, hatred, or delusion, then we keep creating formations that support a trend or habit. This can create a mental movement towards or away from an object. There may be confusion and ignorance regarding the activities that we're producing through body, speech, and mind. Those activities, saṅkhāras, further manifest in ways that reinforce old patterns of perception and the accompanying feeling tone.

The feeling tone of unpleasant, pleasant, or neutral, conditions the next uprising of intentional activity, creating more saṅkhāras. The flow keeps going. This endless motion keeps propelling us through time and experience. It's rare that we allow things to slow down and stop.

Another word that we've been talking about recently is \emph{papañca}. In a discussion with Anuruddha, the Buddha lists eight thoughts of a great being (AN 8.30). In this teaching, the Buddha says, ``This Dhamma is for one who delights in \emph{nippapañca}, not for one who delights in \emph{papañca}.''

There are various ways to translate the word ``\emph{papañca}'', including objectification, proliferation, or complication. We engage in papañca with the sense of I, me, or mine: the sense of self being at the base of it.

Papañca can be seen as the emperor of \emph{vaci-saṅkhāra}, verbal formations or verbal conditioners. Papañca is like vaci-saṅkhāra blown up on steroids and fueled by greed, hatred, and delusion. When there is a strong \emph{kilesa} (afflictive state) in the mind and it engages with the thinking process, it causes a runaway-train effect. It accelerates, compounds, and feeds on itself. It becomes its own beast that then turns back on us.

We become the victim of this beast that has been created through clinging and proliferation. What starts out as a perception and feeling based on contact with the outside world is complicated by thinking and proliferation. The process is reinforced over and over again.

With my background as a registered nurse, I see cancer as a metaphor to this process of proliferation. Cancer starts with one aberrant cell that gets fed and nurtured and is not recognized and destroyed by the immune system. It's out of the range of detection. It starts to reproduce and grow and proliferate. In the same way, our minds do this when there's no mindfulness. When there's no oversight, mindfulness isn't present to pick up on the little errant thought that arises out of aversion, craving, or confusion. We feed it and it starts to increase and proliferate. It takes over, like when cancer starts locally and then metastasizes to other parts of the body. The flow of thoughts, feelings, and emotions suddenly goes on a tangent and lands and grows on something else: another memory, injustice, or form of craving. The whole process is self-reinforcing. We've got a fulminating disease, as they would say in the medical world.

On the mental level, these kinds of uncontrollable processes are saṅkhāras, papañca as the extreme manifestation. We are trying to get a handle on and see them so we can learn not to be the helpless victim. The Noble Eightfold Path provides the means for curing the disease.

Through sustained practice, we can understand how to let that process slow down and come to stillness. In the treatment of cancer, sometimes attempts to annihilate and control the cancer cells are effective. Sometimes the effect is only temporary because we haven't completely eliminated the cancer cells. Although we've suppressed them for a while, they come back. Similarly, we can try to annihilate our outward-flowing papañca and place strong holds on it. This may work for periods of time, but we're really trying to get down to the root causes. We use \emph{sati} \emph{sampajañña}, mindfulness and clear comprehension, to stimulate our mental immune system in order to recognize and turn away from the causes that start the whole process.

Part of rooting out the causes is learning how to slow down---to recognize and let go of the underlying tendencies. We try to stop reacting to, feeding, enabling, and encouraging them. It's a process of release, of recognizing and letting go. ``Knowing and letting go'' as Ajahn Chah would say. We let go not out of resistance or aversion or pushing away, but because we see that it's painful to hold onto something that's hot and burning.

If we put our hands on a hot burning object, we immediately release it. It is a reflex. Similarly, we learn to recognize the habits which burn us when we grasp onto them, and we let them go. The more we practice, the quicker we let go. We recognize the true nature of formations, where they come from and how they pass away. Aniccā vata saṅkhārā. All compounded things are impermanent.

Thinking of saṅkhāra as a flow rather than an object gives a sense of movement. Saṅkhāra is something that's in motion and is inconstant. However, when there's an experience of pleasure, we want to grasp on to it. We want to pin it down, solidify it, and keep it. We do that even with experiences that are uncomfortable.

Grasping is based on a perception of ``me,'' a self that is trying to hold on and protect. The experience gels, hardens, ``firms up,'' as Ajahn Sucitto says, into what feels like a solid reality or a solid perception. At the time, we want to hold on to the familiar, even if it's unpleasant, difficult, or harmful because we seek stability. But if we see saṅkhāra as constantly flowing and changing no matter what we try to do, we eventually realize the futility of trying to hold onto it or push it away. We relax into letting go. We allow what is always changing to do just that, according to its own nature.

This is something that we see more clearly with a still mind. If we don't jump on the moving train and go with all the sights, sounds, smells, tastes, tactile sensations, and thoughts, then we are more with the natural flow of conditions. We stand beside the track and watch the train pass by.

``\emph{Uppajjitvā nirujjhanti}'' is another line from the aforementioned reflection, and is translated in different ways. Ajahn Geoff (Ṭhānissaro Bhikkhu) translates it as, ``They (fabrications or compounded things) disband as they are arising.'' It points to the momentary nature of all experience. When the mind is quiet and steps back to observe the flow, arising and passing away are not timebound. Passing away is inherent with arising. Arising and passing away are part of the same experience and to see this deeply is liberating because the natural impulse is to let go.

If all conditions are constantly changing, there's no need to create all the formations that help support the illusion of solidity. When we can do that, in that stilling there is peace. The flow of conditions still goes on, but the mind steps back. There is just the knowing of the flow of experience, the basic function of citta: the quality of awareness.

How can we protect and maintain some of that stillness as we turn back into daily activity? We start with the basics: keeping our precepts.

If we want to have a quieter mind, without the regrets and disturbances that arise from harmful or agitating actions, following moral precepts is our foundation. It's our front line of protection. We cherish the life of even the smallest of creatures, such as mosquitoes or garden pests. We cherish life enough so that although we may try to get them to move away or protect ourselves from them, we don't annihilate them. Reflecting on gentleness and caring for sentient beings creates ease in the mind, a sense of quiet and stillness through non-harming.

It's the same with all the other precepts. The second precept is to not take that which is not given to us, to not steal anything. We don't have to suffer with the remorse of that action, even on small levels. For example, we might think nobody's going to notice some minor form of theft that ``everybody does,'' like cheating on taxes. However, we don't want to have that niggling in the back of the mind.

The third precept is refraining from sexual misconduct, not taking advantage of others in a sensual way. The fourth precept is refraining from lying. Breaking these precepts will bring agitation and discomfort to the mind. Being careful to protect truthfulness and other wholesome qualities of speech is important to our peace of mind. We refrain from harsh speech, backbiting and even \emph{samphappalāpa}, nattering and purposeless chatter. When we follow these wholesome courses of action, we experience a greater level of settledness and stillness.

We can also find ways to simplify our lives. It takes more effort outside of the monastery because not everybody else is doing it. We can question how much we need and ask ourselves where the acquisitive mind takes us. The world is getting more and more complicated with all the various forms of stimulation and distraction that seem to be shortening our attention span, day by day. It's incredible. In waiting rooms, airports, and even walking down the street, people are totally absorbed in their electronics, listening to something with their earbuds. We can reflect on how this affects us and perhaps learn to not shut off from experiencing the world around us in self-absorption.

Take time to ask yourself, ``What am I doing? Do I need to do this? How's this affecting my mind? How's this affecting the quality of my attention? Is it leading to inner stillness? Is it supporting mindfulness and clear comprehension?'' Take a break. Step back and remind yourself of the state of mind you want.

It's easy to get involved in work, activity, and stimulation. Doing things, picking up the pace and saying, ``Okay, in three months' time I'll get to go on retreat again.'' We go back and forth between living an intensely stimulating, distracting life and shutting it all off with special conditions, even coming to the monastery for a period of time. That's good to do. It's better than not doing it. But try to be consistent with the practice in daily life and really question the level of distraction you engage in.

If you're in the car as a passenger, put down the phone for half an hour and just sit quietly. Try to do that. Take a break.

If we're intent on preserving the quality of stillness, work with sense restraint, protecting how we attend to the outside world. We do have to engage with and experience sights, sounds, and tactile sensations, but \emph{how} we engage is up to us. Without putting ourselves in sensory deprivation tanks and walling ourselves off from every possible kind of stimulation, we can watch how we attend so that we don't get fooled by the exaggerated qualities of the world around us. Whether it's the exaggerated attractions or aversions, we don't have to buy into the mask, the illusion, the enchantment of sights, sounds, and objects that catch our attention. Instead, we look at them with wisdom, in their true nature.

We can absorb into objects that we find attractive because we attend to signs and features that stimulate and support desire and craving. We can counteract that absorption by looking at things more completely, looking at things in their context and getting a true picture of what we find so attractive. Is it really inherently attractive? What is it that seems to be attractive? What is desire going out to look for? Notice that process and give it a reality check with wise attention.

Desire exaggerates the beautiful qualities of an object. It holds on to the object and builds on it. We can stop and say, ``Oh, I am creating a reality that's not really there.'' It's the same with aversion and dislike, how we create and compound stories, one on top of another, through inappropriate attention. We attend to the undesirable qualities and try to get away from them. We create a whole world around this perception. We can think, ``Maybe there are some good qualities of that person whom I find irritating or difficult.'' Try to see the whole person.

Remember the human qualities we all have, including the undesirable or irritating qualities. We can realize, ``Okay, they're not circumscribed to one being who happens to be in my field of attention right now. These are qualities that arise and pass away, too.'' We don't have to hold on to them or run away from them. But we have to protect our minds from going out and getting lost in them.

We see the danger in lack of sense restraint. We step back and see the flow of the world, and we don't jump into it. We try not to get entangled. We can't stop the flow of \emph{saṃsāra}, but we can realize peace by dropping back and resting in the knowing mind, in awareness. Maybe we've had a few momentary glimpses of peace from our dhamma practice that give us a little bit of insight and inspiration: ``Maybe what the Buddha says is true. Maybe there is a path that can go to the complete cessation of entanglement.''

Upāsikā Kee Nanayon described this kind of knowing. ``An inward-staying, unentangled knowing, all outward-going knowing cast aside.''\footnote{Upāsikā Kee Nanayon, \emph{An Unentangled Knowing: The Teachings of a Thai Buddhist Lay Woman}, trans. Thānissaro Bhikkhu (Valley Center, CA: Dhamma Dana Publications, 1995), 29, dhammatalks.org.} Our awareness moves inward: knowing, watching, letting go, and not jumping into the flow.

In that complete stilling is true peace.
